\documentclass[11pt]{article}
\usepackage[margin=1in]{geometry}
\title{SNc, Lesioned Hemisphere --- Interpretation}
\author{GSE233866, 6-OHDA lesion cohort}
\date{}

\begin{document}
\maketitle

6-OHDA lesioning collapses the SNc population from 8{,}637 to 2{,}291 cells
(3.8$\times$ loss), and the network built from what remains needs a higher
soft power (7, $R^2=0.905$) than the intact-hemisphere network (6) to reach
the same scale-free fit quality --- expected behavior when a smaller,
survivor-biased cell population produces a noisier correlation structure.

CACNA1D moves to the \textbf{brown} module with $\mathrm{kME}=0.398$, up from
0.346 in the intact network. Read qualitatively, this points the same
direction as the human PD comparison elsewhere in this project: CACNA1D
becomes \emph{more} central to its co-expression module after
disease/lesioning, not less. This standalone comparison is suggestive but not
a statistical test in itself --- the two networks were clustered
independently, so module identity and kME are not directly comparable at
the significance level. The formal version of this claim is the DME test in
\texttt{lesioned\_vs\_intact\_combined\_dme\_snc/}, which pools both
conditions into one shared network before testing.

\end{document}
